% Вставляется в main.tex до leds.tex
\section{Язык Lua}
\subsection*{Цели раздела}
\begin{itemize}
  \item Освоить типы данных, области видимости и таблицы.
  \item Научиться работать с арифметикой, логикой, условиями и циклами.
  \item Использовать идиомы Lua для краткости и надёжности кода.
  \item Подготовить основу для программирования индикации и миссий.
\end{itemize}
\subsection{Введение}
Lua — лёгкий и встраиваемый язык программирования, созданный в начале 1990‑х для расширения приложений скриптами. Он компактен, прост и быстро работает, поэтому широко используется в встраиваемых системах, играх и робототехнике.

\subsection{Толковый словарь}
Нажмите на термин в тексте, чтобы перейти к его определению.
\begin{description}
  \item[\hypertarget{term:number}{Число}] вещественное значение (обычно двойной точности), например \verb|3.14| или \verb|10|; доступна математика через \verb|math|.
  \item[\hypertarget{term:string}{Строка}] последовательность символов в кавычках, например \verb|"Lua"|; операции: конкатенация \verb|..|, длина \verb|#s|.
  \item[\hypertarget{term:boolean}{Булево}] логическое значение \verb|true| или \verb|false|; используется в условиях.
  \item[\hypertarget{term:nil}{nil}] «пусто»/«нет значения»; удаляет поля таблицы, означает отсутствие данных.
  \item[\hypertarget{term:table}{Таблица}] универсальная структура (массив и словарь одновременно); хранит позиции \verb|t[1]| и поля \verb|t.x|.
  \item[\hypertarget{term:function}{Функция}] блок кода, вызываемый по имени, принимает аргументы и возвращает значения.
  \item[\hypertarget{term:type}{type()}] функция, возвращающая тип значения: \verb|type(10) -> "number"|.
  \item[\hypertarget{term:scope}{Область видимости}] \verb|local| делает переменную видимой только в текущем блоке; без \verb|local| — глобальная.
  \item[\hypertarget{term:operator}{Оператор}] символ или слово, выполняющее действие: \verb|+ - * / % ^ and or not|.
  \item[\hypertarget{term:expression}{Выражение}] комбинация значений и операторов, дающая новое значение.
  \item[\hypertarget{term:condition}{Условие}] логическое выражение, определяющее выбор ветки выполнения.
  \item[\hypertarget{term:loop}{Цикл}] повторение действий: \verb|for|, \verb|while|, \verb|repeat until|.
  \item[\hypertarget{term:ipairs}{ipairs/pairs}] итераторы по таблицам: \verb|ipairs| для массивов; \verb|pairs| для словарей (порядок не гарантирован).
  \item[\hypertarget{term:truthiness}{Истинность}] в Lua ложными считаются только \verb|false| и \verb|nil|.
  \item[\hypertarget{term:unpack}{table.unpack}] распаковывает таблицу в список аргументов: \verb|table.unpack({1,0,0}) -> 1,0,0|. Это функция из библиотеки \verb|table|, а не глобальная функция: вызывать её нужно как \verb|table.unpack|. Для краткости в начале скрипта часто пишут \verb|local unpack = table.unpack|.
\end{description}

\subsection{История и назначение}
- Появился как средство «подключаемой логики» поверх приложений на C/C++.
- Минимальные требования к памяти, переносимость и простой синтаксис.
- Табличные структуры данных (\verb|table|) и функции первого класса делают язык выразительным при небольшом ядре.

\subsection{Lua в нашем курсе}
Ниже — краткий анонс того, для чего нам понадобится Lua; каждый из этих инструментов подробно разбирается в своей главе далее, здесь их не нужно запоминать.
- Мы используем Lua, чтобы программировать поведение квадрокоптера «Геоскан Пионер».
- Через модуль \verb|Ledbar| управляем светодиодной индикацией.
- Через \verb|Timer| создаём неблокирующие анимации.
- Через \verb|callback(event)| реагируем на системные события (посадка, напряжение, достижение точки).

\subsection{Базовые конструкции Lua}
Ниже — краткий пример базовых элементов: переменная, строка и таблица. Циклы и условия разбираются далее.
\begin{infobox}[Практика]
    Этот и последующие примеры кода вы можете копировать в редактор Pioneer Station и запускать, чтобы увидеть результат в консоли.
\end{infobox}
\begin{codefile}{lua/examples/04_language/lua_intro_basics.lua}
\end{codefile}
\subsection*{Итоги раздела}
Вы научились объявлять и различать типы данных, писать выражения, условия и циклы, работать с таблицами как с массивами и словарями, применять стандартные математические функции. Эти навыки будут использованы в главах о светодиодах и автономных миссиях.
